%% ============================================================
%%  ejecucion.tex  —  Sección 1.6
%% ============================================================

\subsection{Ejecución de experimentos}

\subsubsection*{Experimentos 1–3: benchmark unificado (\texttt{main\_benchmark.cpp})}

Los Experimentos 1–3 se ejecutaron con semilla fija $42$. Los resultados se
almacenaron en \texttt{aporte\_propio/results.csv} (16 mediciones).

\paragraph{Experimento 1 — Escalabilidad en $n$ (disperso, $m=2n$)}

La Tabla~\ref{tab:exp1} muestra los tiempos de ejecución en $\mu$s para
grafos dispersos con $m \approx 2n$.
Bellman-Ford se omite (\textsc{skip}) para $n > 2\,000$ por su coste
$O(m\cdot n)$.

\begin{table*}[t]
  \centering
  \caption{Tiempos ($\mu$s) y operaciones — Experimento 1 ($m=2n$, semilla 42).}
  \label{tab:exp1}
  \small
  \begin{tabular}{rrrrrrr}
    \toprule
    \textbf{$n$} & \textbf{$m$} &
    \textbf{Dijkstra} & \textbf{BF} &
    \textbf{DMMSY} & \textbf{Thorup} &
    \textbf{Ops (Dij)} \\
    \midrule
    100    & 200     & 10    & 11   & 106   & 960    & 200   \\
    500    & 1\,000  & 50    & 52   & 438   & 3\,127 & 1\,000 \\
    1\,000 & 2\,000  & 129   & 102  & 1\,108 & 5\,779 & 1\,994 \\
    2\,000 & 4\,000  & 244   & 251  & 1\,232 & 10\,877 & 4\,000 \\
    5\,000 & 10\,000 & 581   & \textsc{skip} & 3\,198 & 25\,742 & 9\,995 \\
    10\,000 & 20\,000 & 1\,380 & \textsc{skip} & 7\,492 & 52\,209 & 20\,000 \\
    \bottomrule
  \end{tabular}
\end{table*}

\noindent\textbf{Observación clave:} De $n=100$ a $n=10\,000$ (factor
$100\times$ en $n$), el tiempo de Dijkstra creció $138\times$, consistente
con $O(n\log n)$ para $m=O(n)$ (factor teórico esperado: $\approx 200\times$).
Thorup mostró el mayor crecimiento absoluto ($\approx 54\times$) con un
overhead constante de $\approx 35$--$40\times$ sobre Dijkstra.

\paragraph{Experimento 2 — Escalabilidad en $m$ ($n=1\,000$ fijo)}

\begin{table*}[t]
  \centering
  \caption{Tiempos ($\mu$s) — Experimento 2 ($n=1\,000$, semilla 42).}
  \label{tab:exp2}
  \small
  \begin{tabular}{rrrrrr}
    \toprule
    \textbf{$m$} & \textbf{Densidad} &
    \textbf{Dijkstra} & \textbf{BF} & \textbf{DMMSY} & \textbf{Thorup} \\
    \midrule
    1\,000   & 0.10\% & 53  & 1\,315 & 692 & 5\,334 \\
    5\,000   & 0.50\% & 301 & 200    & 740 & 7\,152 \\
    10\,000  & 1.0\%  & 278 & 223    & 599 & 5\,208 \\
    50\,000  & 5.0\%  & 631 & 559    & 658 & 4\,878 \\
    100\,000 & 10.0\% & 619 & 857    & 793 & 5\,492 \\
    200\,000 & 20.0\% & 838 & 1\,439 & 883 & 6\,245 \\
    \bottomrule
  \end{tabular}
\end{table*}

\noindent\textbf{Observación:} El punto de cruce BF $\approx$ Dijkstra
ocurre entre $m=5\,000$ y $m=10\,000$. Thorup presenta un comportamiento
casi constante ($\approx 5\,000$--$7\,000\,\mu$s) independiente de $m$,
coherente con el dominio del término $O(\sqrt{w_{\max}\cdot n}) \approx 3\,162$
sobre el término lineal en $m$ para $n=1\,000$.

\paragraph{Experimento 3 — Densidades ($n=500$ fijo)}

\begin{table*}[t]
  \centering
  \caption{Tiempos ($\mu$s) — Experimento 3 ($n=500$, semilla 42).}
  \label{tab:exp3}
  \small
  \begin{tabular}{rrrrrr}
    \toprule
    \textbf{$m$} & \textbf{Etiqueta} &
    \textbf{Dijkstra} & \textbf{BF} & \textbf{DMMSY} & \textbf{Thorup} \\
    \midrule
    500     & muy disperso & 15  & 366 & 1\,415 & 3\,218 \\
    2\,500  & disperso     & 383 & 231 & 588    & 3\,367 \\
    12\,500 & medio        & 348 & 157 & 470    & 2\,807 \\
    50\,000 & denso        & 683 & 570 & 533    & 2\,968 \\
    \bottomrule
  \end{tabular}
\end{table*}

\noindent\textbf{Observación:} Para $m=50\,000$ ($\text{densidad}=20\%$),
DMMSY ($533\,\mu$s) supera a Dijkstra ($683\,\mu$s): es el único escenario
de los Experimentos 1–3 donde la implementación DMMSY resulta más rápida.

\subsubsection*{Experimento 4 — Aporte propio: impacto de la densidad}

El Experimento 4 amplía el Experimento 3 a un diseño factorial completo con
$n \in \{500, 1\,000, 2\,000, 5\,000\}$, 5 niveles de densidad y 3 semillas.

\paragraph{Resultado principal: tiempo vs.\ densidad ($n=1\,000$)}

La Figura~\ref{fig:g1} muestra los tiempos promedio con barras de error
$\pm 1\sigma$ para $n=1\,000$.

\begin{figure}[H]
  \centering
  \includegraphics[width=0.92\linewidth]{grafica1_tiempo_vs_densidad.png}
  \caption{Tiempo de ejecución vs.\ densidad del grafo para $n=1\,000$.
    Barras de error: $\pm 1$ desviación estándar sobre 3 semillas.
    Bellman-Ford muestra la inversión característica: más lento que Dijkstra
    en grafos muy dispersos, más rápido en densidades intermedias.}
  \label{fig:g1}
\end{figure}

La Tabla~\ref{tab:exp4_n1000} detalla los valores numéricos para $n=1\,000$.

\begin{table*}[t]
  \centering
  \caption{Tiempos promedio ($\mu$s $\pm$ desv.\ est.) — Exp.\ 4, $n=1\,000$,
    3 semillas.}
  \label{tab:exp4_n1000}
  \scriptsize
  \begin{tabular}{lrrrrl}
    \toprule
    \textbf{Densidad} & \textbf{Dijkstra} & \textbf{BF} &
    \textbf{Thorup} & \textbf{DMMSY} & \textbf{Ganador} \\
    \midrule
    muy\_disperso & $38.0 \pm 7.1$    & $1387.7 \pm 529.1$ & $13712.0 \pm 676.7$ & $1415.0 \pm 230.8$ & Dijkstra \\
    disperso      & $202.7 \pm 56.3$  & $157.3 \pm 50.5$   & $6912.3 \pm 1419.2$ & $754.7 \pm 215.7$  & Bellman-Ford \\
    semi\_denso   & $215.0 \pm 43.2$  & $147.0 \pm 15.6$   & $6344.0 \pm 631.5$  & $636.7 \pm 49.0$   & Bellman-Ford \\
    denso         & $310.7 \pm 32.7$  & $217.3 \pm 10.9$   & $6001.0 \pm 411.1$  & $669.3 \pm 84.6$   & Bellman-Ford \\
    muy\_denso    & $1220.3 \pm 110.5$ & $2741.7 \pm 1112.8$ & $7060.0 \pm 600.8$ & $1858.0 \pm 542.0$ & Dijkstra \\
    \bottomrule
  \end{tabular}
\end{table*}

\paragraph{Operaciones de relajación}

La Figura~\ref{fig:g2} muestra las operaciones realizadas por cada algoritmo
en función de $m$, para $n \in \{500, 1\,000, 2\,000\}$.

\begin{figure}[H]
  \centering
  \includegraphics[width=\linewidth]{grafica2_ops_vs_m.png}
  \caption{Operaciones de relajación vs.\ $m$ (escala log-log).
    Dijkstra, Thorup y DMMSY realizan exactamente el mismo número de
    operaciones (curvas superpuestas). Bellman-Ford (eje derecho, naranja)
    realiza entre $8\times$ y $163\times$ más relajaciones.}
  \label{fig:g2}
\end{figure}

\paragraph{Mapa de ganadores por configuración}

La Figura~\ref{fig:g3} representa el algoritmo más rápido en cada celda
$(n, \text{densidad})$ del espacio experimental.

\begin{figure}[H]
  \centering
  \includegraphics[width=0.88\linewidth]{grafica3_heatmap_ganador.png}
  \caption{Heatmap del algoritmo ganador (menor tiempo promedio) por
    configuración. Azul: Dijkstra; Naranja: Bellman-Ford; Rojo: DMMSY.
    Thorup no resultó ganador en ninguna celda.}
  \label{fig:g3}
\end{figure}

\noindent\textbf{Observaciones del mapa:}
\begin{itemize}
  \item \textbf{Dijkstra} domina la columna \texttt{muy\_disperso} en todos los
    valores de $n$, y la fila completa para densidades muy altas ($n \leq 1\,000$,
    \texttt{muy\_denso}).
  \item \textbf{Bellman-Ford} gana en la zona media (\texttt{disperso} a
    \texttt{denso}) para $n \leq 2\,000$, gracias a la terminación temprana en
    grafos aleatorios dirigidos sin ciclos negativos.
  \item \textbf{DMMSY} supera a Dijkstra en la esquina de alta densidad y
    $n \geq 2\,000$: para $n=5\,000$, \texttt{muy\_denso}, DMMSY tarda
    $10\,171\,\mu$s frente a $11\,992\,\mu$s de Dijkstra ($\approx 15\%$
    de mejora).
\end{itemize}

\paragraph{Escalabilidad con $n$ — densidad dispersa ($m = 2n$)}

\begin{figure}[H]
  \centering
  \includegraphics[width=0.88\linewidth]{grafica4_scaling_n.png}
  \caption{Escalabilidad con $n$ para $m=2n$ (disperso).
    Líneas punteadas grises: referencias teóricas $O(n\log n)$ y $O(n^2)$
    normalizadas al punto $n=500$ de Dijkstra.
    El crecimiento empírico de Dijkstra ($4.86\times$ de $n=500$ a $n=5\,000$)
    es sublineal respecto al factor teórico $O(n\log n) \approx 13.7\times$,
    lo que refleja la eficiencia de la caché para $m$ pequeño.}
  \label{fig:g4}
\end{figure}

\begin{table*}[t]
  \centering
  \caption{Escalabilidad con $n$ — densidad \texttt{disperso} ($m=2n$),
    promedio 3 semillas.}
  \label{tab:scaling}
  \small
  \begin{tabular}{rrrrrr}
    \toprule
    \textbf{$n$} & \textbf{$m$} &
    \textbf{Dijkstra ($\mu$s)} & \textbf{BF ($\mu$s)} &
    \textbf{Thorup ($\mu$s)} & \textbf{DMMSY ($\mu$s)} \\
    \midrule
    500   & 1\,000  & 154.7  & 117.3  & 7\,759.0  & 959.3  \\
    1\,000 & 2\,000 & 202.7  & 157.3  & 6\,912.3  & 754.7  \\
    2\,000 & 4\,000 & 431.3  & 311.3  & 12\,236.3 & 1\,639.3 \\
    5\,000 & 10\,000 & 752.3 & ---    & 25\,858.3 & 3\,219.7 \\
    \bottomrule
  \end{tabular}
\end{table*}

\paragraph{Verificación de hipótesis}

\begin{description}
  \item[\textbf{H1 — CONFIRMADA.}] ~\\
    Dijkstra es el ganador en grafos muy dispersos para todos los valores de
    $n$ evaluados (e.g., $38.0\,\mu$s para $n=1\,000$, \texttt{muy\_disperso}).
  \item[\textbf{H2 — PARCIALMENTE CONFIRMADA.}] ~\\
    Se confirma que BF es inviable para $n>2\,000$. Sin embargo, para
    $n \leq 2\,000$ y densidades intermedias, BF \textit{supera} a Dijkstra
    (e.g., $157.3\,\mu$s vs.\ $202.7\,\mu$s para $n=1\,000$, \texttt{disperso}).
    La hipótesis subestimó la eficiencia de la terminación temprana.
  \item[\textbf{H3 — CONFIRMADA.}] ~\\
    Thorup es el algoritmo más lento en todos los escenarios (mínimo:
    $6\,001\,\mu$s para $n=1\,000$, \texttt{denso}).
  \item[\textbf{H4 — CONFIRMADA.}] ~\\
    DMMSY supera a Dijkstra en grafos de muy alta densidad y $n \geq 2\,000$
    ($2\,621\,\mu$s vs.\ $3\,026\,\mu$s para $n=2\,000$, \texttt{muy\_denso}).
\end{description}
